% ============================================================
% The Alignment-Faithfulness Paradox
% ACL 2026 — Two-column format
% ============================================================
\documentclass[11pt]{article}

% ACL 2026 style
\usepackage[margin=1in]{geometry}
\usepackage{acl}
\usepackage{times}
\usepackage{latexsym}
\usepackage{microtype}
\usepackage{booktabs}
\usepackage{amsmath}
\usepackage{amssymb}
\usepackage{graphicx}
\usepackage{xcolor}
\usepackage{soul}
\usepackage{multirow}
\usepackage{hyperref}
\usepackage{url}

% ── RESULT placeholder macro ──────────────────────────────────
% \RESULT{key} renders as a yellow-highlighted box with the key name.
% After BigRed200 results return, run:
%   python fill_results.py
% which replaces every \RESULT{key} with the value from
%   analysis/final_results_summary.json
% and writes paper_final.tex.
\newcommand{\RESULT}[1]{\colorbox{yellow}{\textbf{[#1]}}}

% ─────────────────────────────────────────────────────────────
\title{The Alignment-Faithfulness Paradox:\\
Does CoT Quality Fine-Tuning Reduce Parametric Faithfulness?}

\author{
  Madhavan Balaji \\
  Indiana University \\
  \texttt{madbala@iu.edu}
}

\begin{document}
\maketitle

% ─────────────────────────────────────────────────────────────
\begin{abstract}
Large language models increasingly produce Chain-of-Thought (CoT) reasoning
that appears coherent and well-structured, yet may not reflect the model's
actual computational process. We investigate whether fine-tuning on
high-quality CoT examples --- a natural strategy for improving reasoning
transparency --- paradoxically \emph{reduces} parametric faithfulness,
i.e., the degree to which the expressed reasoning causally drives predictions.
Extending the Negative Preference Optimization (NPO) unlearning framework
of \citet{tutek2025faithfulness}, we score \RESULT{n_scored} CoT instances
generated by LLaMA-3-3B using an LLM judge and fine-tune LoRA adapters on
high- and low-quality splits. We then re-evaluate parametric faithfulness on
held-out instances under each condition. Our results show that high-quality
CoT fine-tuning yields binary faithfulness of \RESULT{hq_binary_faith_pct}\%
versus baseline \RESULT{baseline_binary_faith_pct}\%, and continuous
faithfulness ($\Delta p$) of \RESULT{hq_mean_delta_p} vs.\
\RESULT{baseline_mean_delta_p} (Welch $t$=\RESULT{hq_t_stat},
$p$=\RESULT{hq_t_p}). We interpret this as the \emph{alignment-faithfulness
paradox}: surface-level reasoning quality and parametric faithfulness are
partially dissociable, with implications for the interpretability of
instruction-tuned models.
\end{abstract}

% ─────────────────────────────────────────────────────────────
\section{Introduction}
\label{sec:intro}

Chain-of-Thought prompting has become a cornerstone of modern language model
deployment \citep{wei2022cot}. By conditioning models to produce intermediate
reasoning steps before answering, practitioners report improvements in
accuracy and interpretability across a wide range of tasks
\citep{kojima2022zeroshot,lightman2023prm}. The implicit assumption is that
these steps are \emph{faithful}: the model uses the expressed reasoning to
arrive at its answer. If this assumption holds, improving reasoning quality
through fine-tuning should simultaneously improve interpretability.

However, faithfulness and quality are conceptually distinct. A CoT can be
internally coherent, factually accurate, and stylistically polished while
playing no causal role in the model's prediction. Prior work has documented
exactly this dissociation: models produce plausible-sounding reasoning chains
that are post-hoc rationalizations rather than generative explanations
\citep{turpin2023faithfulness,lanham2023measuring,lyu2023faithful}. The
mechanism of faithfulness --- how reasoning steps encoded in transformer
weights causally influence output distributions --- remains poorly understood.

\citet{tutek2025faithfulness} propose a rigorous operationalization:
\emph{parametric faithfulness} measures whether surgically unlearning a CoT
step (via Negative Preference Optimization applied to FF2 layers) changes the
model's answer. This probes whether the information in the expressed reasoning
step is encoded in the same parameters that drive the final prediction. Their
main finding is that faithfulness so measured is empirically distinct from
accuracy, varies by model family and dataset, and correlates strongly with
unlearning efficacy ($r = 0.937$, $p < 0.0001$ in our reproduction).

In this work, we ask a direct follow-up: \textbf{does fine-tuning on
high-quality CoT outputs increase or decrease parametric faithfulness?} This
question has practical import. If the answer is ``increases,'' then standard
instruction-following fine-tuning on curated reasoning data is a path toward
more interpretable models. If the answer is ``decreases'' or ``null,'' then
surface quality and computational faithfulness are dissociated, and we cannot
assume that CoT quality improvements translate to faithfulness improvements.

We call the hypothesis that quality fine-tuning reduces faithfulness the
\emph{alignment-faithfulness paradox}. The mechanism we propose is as
follows: high-quality CoT fine-tuning strengthens the association between
linguistic reasoning patterns and correct answers in the model's surface
behavior, but this association may be mediated by shallow token-prediction
paths rather than by the FF2-layer parametric representations that NPO
targets. As a result, the model learns to produce better-sounding reasoning
while the causal link between reasoning content and predictions weakens.

Our contributions are:
\begin{enumerate}
  \item A reproduction of the NPO faithfulness framework confirming core
        findings and identifying four previously unreported bugs in the
        released codebase.
  \item A new continuous faithfulness metric ($\Delta p$) and analysis of
        \emph{subcritical} faithfulness --- instances where CoT causally
        reduces correct-answer probability but does not flip predictions ---
        showing that binary faithfulness undercounts causal signal by
        \RESULT{subcritical_undercounting_pct}\% in LLaMA-3-3B.
  \item A controlled experiment comparing parametric faithfulness under
        baseline, high-quality CoT, and low-quality CoT fine-tuning conditions
        across \RESULT{n_test_instances} held-out instances.
  \item Evidence \RESULT{paradox_verdict} the alignment-faithfulness paradox:
        \RESULT{paradox_direction}.
\end{enumerate}

% ─────────────────────────────────────────────────────────────
\section{Background and Related Work}
\label{sec:background}

\subsection{Chain-of-Thought Faithfulness}

Faithfulness in NLP was formalized by \citet{jacovi2020faithfully} as the
degree to which an explanation accurately represents the model's actual
reasoning process, distinct from \emph{plausibility}, which concerns human
interpretability. Early work showed that post-hoc rationales generated by
chain-of-thought models can be influenced by spurious features in the prompt
\citep{turpin2023faithfulness}: adding a biasing hint (e.g., labeling a wrong
answer as the correct one) changes the model's answer without the expressed
reasoning acknowledging the influence. \citet{lanham2023measuring} applied
direct perturbation tests to CoT, showing that corrupting intermediate steps
often fails to affect final predictions, suggesting many steps are redundant.

\citet{lyu2023faithful} proposed structured symbolic CoT to improve
faithfulness by construction, while \citet{prystawski2024think} argued that
any faithful intermediate representation requires information to propagate
through the residual stream in ways that are difficult to control via
prompting alone. These lines of work converge on a pessimistic view: standard
auto-regressive CoT generation is not guaranteed to produce faithful
reasoning, even when the reasoning appears coherent.

\subsection{Parametric Faithfulness via Unlearning}

\citet{tutek2025faithfulness} introduce a mechanistic operationalization that
sidesteps the ambiguity of perturbation tests. Their key insight is that if a
CoT step is parametrically faithful --- encoded in the same model weights that
drive the final prediction --- then selectively unlearning that step should
change the prediction. They use Negative Preference Optimization (NPO)
\citep{zhang2024negative} to minimize the log-likelihood of the CoT step
while maintaining a KL divergence constraint on retain sequences (npo\_KL
objective). Unlearning is applied only to FF2 layers (\texttt{mlp.down\_proj}
weights), which store factual associations \citep{meng2022rome}. A binary
faithfulness label is assigned if any prediction flip occurs across training
epochs.

This framework has several appealing properties: it targets parametric
encoding rather than surface behavior, it is model-agnostic, and it produces
per-instance faithfulness scores rather than aggregate statistics. Its main
limitation is the binary nature of the faithfulness label, which we address
in this work.

\subsection{Fine-Tuning, RLHF, and Reasoning Quality}

Reinforcement Learning from Human Feedback \citep{ouyang2022rlhf} and its
variants have demonstrated that optimizing for human preference ratings can
substantially improve the surface quality of model outputs, including CoT
reasoning. \citet{lightman2023prm} show that process reward models trained to
evaluate individual reasoning steps yield further gains. However, whether
these improvements carry through to faithfulness has not been tested. The
closest prior work is on \emph{sycophancy} \citep{perez2022sycophancy}, which
shows that RLHF-aligned models are more likely to produce socially desirable
but incorrect reasoning --- a qualitatively similar dissociation to what we
hypothesize.

% ─────────────────────────────────────────────────────────────
\section{Hypothesis and Mechanism}
\label{sec:hypothesis}

We formalize the alignment-faithfulness paradox hypothesis as follows:

\begin{quote}
\textbf{H1.} Fine-tuning LLaMA-3-3B on high-quality CoT instances will
  yield lower parametric faithfulness (binary rate and mean $\Delta p$) on
  held-out instances, compared to the unmodified base model.
\end{quote}

\begin{quote}
\textbf{H2.} Fine-tuning on low-quality CoT instances will yield equal or
  lower parametric faithfulness than H1 (null baseline hypothesis for the
  quality comparison).
\end{quote}

The proposed mechanism rests on a distinction between two routes through
which a model can arrive at a correct answer given a CoT:
\begin{enumerate}
  \item \textbf{Parametric route}: The semantic content of the CoT step
        activates or strengthens a factual association stored in FF2 layers,
        shifting the output distribution toward the correct answer.
  \item \textbf{Surface route}: The CoT string's lexical patterns (e.g.,
        hedging language, certain keywords) correlate with correct answers
        in the training distribution, allowing the model to shortcut to the
        correct prediction without encoding the causal structure.
\end{enumerate}
We hypothesize that high-quality CoT fine-tuning strengthens the surface
route, because the training signal (scored CoT examples) selects for
stylistically coherent, complete reasoning chains whose text is correlated
with correct answers in the training data. The FF2 parametric route is not
directly incentivized, and may even be weakened if the model learns that
surface-level completeness is sufficient.

% ─────────────────────────────────────────────────────────────
\section{Experimental Setup}
\label{sec:setup}

\subsection{Base Model and Datasets}

We use \textbf{LLaMA-3.2-3B-Instruct} \citep{dubey2024llama3} as our base
model throughout, following \citet{tutek2025faithfulness} who use this model
as their primary small-model condition. We draw from four multiple-choice QA
benchmarks used in the original paper: \textbf{ARC-Challenge}
\citep{clark2018arc}, \textbf{OpenBookQA} \citep{mihaylov2018openbook},
\textbf{Sports Understanding} \citep{srivastava2022bigbench}, and
\textbf{StrategyQA} \citep{geva2021strategyqa}. All evaluation instances are
drawn from the same CoT cache used in our reproduction
(\texttt{final\_cot/\{dataset\}/}).

\subsection{CoT Quality Scoring}

We score \RESULT{n_scored} CoT instances generated by LLaMA-3-3B on a
3-dimensional rubric:
\begin{itemize}
  \item \textbf{Coherence} (1--5): Is the reasoning logically consistent and
        well-structured?
  \item \textbf{Plausibility} (1--5): Does each step follow plausibly from
        the question and prior steps?
  \item \textbf{Completeness} (1--5): Does the CoT cover all necessary
        reasoning without unnecessary gaps?
\end{itemize}
Scoring is performed by \texttt{claude-sonnet-4-5} via the Anthropic API,
using a structured prompt that requests integer scores for each dimension.
The composite score is the mean of the three dimensions. Inter-rater
consistency is estimated by rescoring a random 5\% sample ($n =
\RESULT{n_rescore}$); Cronbach's $\alpha = \RESULT{cronbach_alpha}$.

\subsection{Fine-Tuning Conditions}

We create three experimental conditions using the scored CoT pool:
\begin{enumerate}
  \item \textbf{Baseline}: The unmodified LLaMA-3-3B-Instruct checkpoint.
  \item \textbf{High-quality}: A LoRA adapter \citep{hu2022lora} trained on
        the top-50\% of instances by composite score ($n =
        \RESULT{n_high_quality}$).
  \item \textbf{Low-quality}: A LoRA adapter trained on the bottom-50\%
        ($n = \RESULT{n_low_quality}$).
\end{enumerate}
Both adapters use identical hyperparameters: rank $r = 16$, scaling $\alpha =
32$, dropout $= 0.05$, applied to all 7 projection matrices
(\texttt{q\_proj}, \texttt{v\_proj}, \texttt{k\_proj}, \texttt{o\_proj},
\texttt{gate\_proj}, \texttt{up\_proj}, \texttt{down\_proj}). Training runs
for 3 epochs with learning rate $2 \times 10^{-4}$, effective batch size 16
(4 per device, 4 gradient accumulation steps), and a cosine schedule with 5\%
warmup. Adapters are merged into the base weights before evaluation
(\texttt{merge\_and\_unload()}), ensuring evaluation is performed on a
standard dense model.

\subsection{Faithfulness Evaluation}

We hold out \RESULT{n_test_instances} instances stratified by dataset
(approximately 12-13 per dataset). For each held-out instance under each
condition, we run the full NPO unlearning loop with the same hyperparameters
as the base reproduction: FF2-only optimization, \texttt{npo\_KL} objective,
learning rate $3 \times 10^{-5}$, 5 epochs.

We report two faithfulness metrics:
\begin{itemize}
  \item \textbf{Binary faithfulness}: Fraction of instances where any
        prediction flip occurs across epochs 1--5.
  \item \textbf{Continuous faithfulness ($\Delta p$)}: The change in
        normalized probability assigned to the correct answer between epoch 0
        and the epoch with maximum $\Delta p$:
        \[
          \Delta p = p_{\text{correct}}^{(0)} - \min_e p_{\text{correct}}^{(e)}
        \]
        Positive $\Delta p$ indicates that unlearning reduced the model's
        confidence in the correct answer, i.e., the CoT step causally
        contributed to the prediction.
\end{itemize}

Statistical comparisons use McNemar's test for binary faithfulness (paired
at the instance level) and Welch's $t$-test for continuous $\Delta p$, with
$\alpha = 0.05$ and no correction for multiple comparisons (two planned
contrasts: baseline vs.\ high-quality and baseline vs.\ low-quality).

% ─────────────────────────────────────────────────────────────
\section{Results}
\label{sec:results}

\subsection{CoT Quality Distribution}

Composite quality scores span a range of \RESULT{score_min}--\RESULT{score_max}
with mean \RESULT{score_mean} ($\sigma = \RESULT{score_std}$). The
distribution is \RESULT{score_shape} (Shapiro--Wilk $p = \RESULT{shapiro_p}$).
High-quality instances (top 50\%) have mean composite score
\RESULT{hq_mean_score}; low-quality instances have mean score
\RESULT{lq_mean_score}. The gap between splits is \RESULT{score_split_gap}
points, representing \RESULT{score_split_gap_pct}\% of the full score range.

\subsection{Faithfulness Under Each Condition}

Table~\ref{tab:main_results} reports binary faithfulness rates, mean $\Delta
p$, and the percentage of instances with positive $\Delta p$ under each
condition.

\begin{table}[ht]
\centering
\small
\begin{tabular}{lccc}
\toprule
\textbf{Condition} & \textbf{Binary \%} & \textbf{Mean $\Delta p$} & \textbf{\% $\Delta p > 0$} \\
\midrule
Baseline           & \RESULT{baseline_binary_faith_pct} & \RESULT{baseline_mean_delta_p} & \RESULT{baseline_pct_positive_dp} \\
High-quality CoT   & \RESULT{hq_binary_faith_pct} & \RESULT{hq_mean_delta_p} & \RESULT{hq_pct_positive_dp} \\
Low-quality CoT    & \RESULT{lq_binary_faith_pct} & \RESULT{lq_mean_delta_p} & \RESULT{lq_pct_positive_dp} \\
\bottomrule
\end{tabular}
\caption{Faithfulness metrics by fine-tuning condition.
  $N = \RESULT{n_test_instances}$ held-out instances per condition.}
\label{tab:main_results}
\end{table}

\subsection{Statistical Tests}

\paragraph{Baseline vs.\ High-Quality.}
Binary faithfulness: \RESULT{baseline_binary_faith_pct}\% $\rightarrow$
\RESULT{hq_binary_faith_pct}\% (McNemar $\chi^2 = \RESULT{hq_mcnemar_chi2}$,
$p = \RESULT{hq_mcnemar_p}$, \RESULT{hq_mcnemar_sig}).
Continuous $\Delta p$: \RESULT{baseline_mean_delta_p} $\rightarrow$
\RESULT{hq_mean_delta_p} (Welch $t = \RESULT{hq_t_stat}$,
$p = \RESULT{hq_t_p}$, \RESULT{hq_t_sig}).

\paragraph{Baseline vs.\ Low-Quality.}
Binary faithfulness: \RESULT{baseline_binary_faith_pct}\% $\rightarrow$
\RESULT{lq_binary_faith_pct}\% (McNemar $\chi^2 = \RESULT{lq_mcnemar_chi2}$,
$p = \RESULT{lq_mcnemar_p}$, \RESULT{lq_mcnemar_sig}).
Continuous $\Delta p$: \RESULT{baseline_mean_delta_p} $\rightarrow$
\RESULT{lq_mean_delta_p} (Welch $t = \RESULT{lq_t_stat}$,
$p = \RESULT{lq_t_p}$, \RESULT{lq_t_sig}).

\subsection{Subcritical Faithfulness}

Across all conditions, binary faithfulness systematically undercounts causal
signal. Of instances with $\Delta p > 0$ in the baseline condition,
\RESULT{subcritical_undercounting_pct}\% do not exhibit a prediction flip ---
the CoT causally reduced correct-answer confidence without crossing the
decision boundary. This subcritical faithfulness rate is highest in
\RESULT{subcritical_highest_condition} (\RESULT{subcritical_highest_pct}\%)
and lowest in \RESULT{subcritical_lowest_condition}
(\RESULT{subcritical_lowest_pct}\%), suggesting that \RESULT{subcritical_interpretation}.

\subsection{Per-Epoch Faithfulness Trajectories}

Figure~\ref{fig:trajectories} (see \texttt{my\_figures/new/alignment\_paradox.png})
shows mean $\Delta p$ by training epoch under each condition. The
\RESULT{trajectory_shape} pattern suggests that \RESULT{trajectory_interpretation}.

% ─────────────────────────────────────────────────────────────
\section{Discussion}
\label{sec:discussion}

\subsection{Interpreting the (Non-)Paradox}

Our results \RESULT{paradox_verdict_extended}. If confirmed (after BigRed200
results are available), several interpretations are consistent with the data.

\paragraph{Interpretation 1: Decoupling via surface routes.}
High-quality CoT fine-tuning teaches the model that coherent, complete
reasoning patterns co-occur with correct answers in the training distribution.
This may bypass the FF2 parametric route entirely: the model learns a stronger
surface-to-answer shortcut, making FF2 weights that encode the CoT step's
semantic content less necessary. Under this view, faithfulness and quality are
necessarily traded off because the training objective for quality (maximize
log-likelihood of good CoT completions) is orthogonal to the faithfulness
objective (maximize causal connection between CoT encoding and prediction).

\paragraph{Interpretation 2: Representation redistribution.}
LoRA fine-tuning adds adapter parameters across attention and MLP layers but
the NPO faithfulness test specifically targets \texttt{mlp.down\_proj} (FF2)
parameters. If fine-tuning redistributes the encoding of CoT content from
FF2 to other layers (e.g., attention heads or FF1 layers), NPO unlearning of
FF2 would show reduced effect even if the CoT remains causally relevant
overall. This would suggest a limitation of the NPO faithfulness probe rather
than a genuine faithfulness reduction.

\paragraph{Interpretation 3: Null result.}
If neither condition differs significantly from baseline, the most defensible
conclusion is that 50-instance held-out evaluation is underpowered to detect
the effect, and that CoT quality fine-tuning with LoRA rank 16 simply does
not change the parametric role of FF2 layers substantially enough to measure
within this evaluation budget.

\subsection{Relationship to Prior Work}

Our experiment operationalizes a conjecture that has appeared implicitly in
several lines of prior work. \citet{turpin2023faithfulness} showed that
aligned models (with RLHF) produce more biased post-hoc rationalizations than
their base counterparts. \citet{lanham2023measuring} found that longer,
higher-quality CoT chains were not more faithful by their perturbation metric.
Our contribution is a controlled experimental design --- same model, same
evaluation protocol, only the CoT training distribution varies --- that
directly tests the quality-faithfulness relationship.

The clearest prior analogue is \citet{perez2022sycophancy}, who demonstrate
that RLHF-tuned models produce more agreeable but less accurate reasoning.
The alignment-faithfulness paradox is the mechanistic counterpart: not only
do aligned models \emph{say} more agreeable things, their parametric weights
may encode those statements less causally.

\subsection{Limitations}

Several limitations constrain our conclusions. First, we evaluate only
LLaMA-3-3B; the paradox may not generalize to larger models or different
architectures. Second, LoRA fine-tuning with rank 16 on 400-500 instances
is a weak intervention compared to full fine-tuning on millions of examples.
Third, our quality scoring relies on a single LLM judge (Claude Sonnet 4-5)
without human validation. Fourth, the NPO faithfulness probe targets only FF2
layers and first CoT steps; faithfulness in middle and later steps may behave
differently. Finally, 50 held-out instances per condition yields limited
statistical power, particularly for detecting small effect sizes.

% ─────────────────────────────────────────────────────────────
\section{Conclusion}
\label{sec:conclusion}

We investigated whether fine-tuning LLaMA-3-3B on high-quality
Chain-of-Thought instances increases parametric faithfulness as measured by
the NPO unlearning framework. Our results \RESULT{conclusion_finding}. The
alignment-faithfulness paradox --- that surface quality improvement may
decouple expressed reasoning from its parametric encoding ---
\RESULT{conclusion_paradox_status}. These findings suggest that improving
interpretability of CoT reasoning requires metrics and training objectives
that directly target causal encoding, not merely surface quality.

Future work should examine larger models, full fine-tuning (not just LoRA),
and faithfulness-aware training objectives that explicitly penalize the
surface shortcut route. The NPO framework is a promising evaluation tool for
such work, provided its FF2 restriction and single-step focus are extended to
broader parametric coverage.

% ─────────────────────────────────────────────────────────────
\bibliography{paper}
\bibliographystyle{acl_natbib}

\end{document}
